\section{Gaussian fourfolds: the polarization, support, and small bound}
\label{sec:gaussian}

\subsection{HD-Q01: the precise principal datum}
Let $\pi:C\to X$ be a connected unramified cyclic cover of degree four, with $g(X)=3$ and generator $\sigma$. Put $D=C/\langle\sigma^2\rangle$. Riemann--Hurwitz gives $g(C)=9$ and $g(D)=5$. The primitive factor is
\[
 B=\operatorname{im}(1-\sigma^2)
  =\ker\bigl(\operatorname{Nm}_{C/D}\bigr)^0\subset J(C).
\]
Here the displayed equality identifies the abelian subvariety $B$ with the indicated connected kernel; the ambient Jacobian has dimension nine. The restriction of its theta class is
\[
 \theta_{J(C)}|_B=2\xi,\qquad
 \xi\ \text{principal},\qquad
 \ker\lambda_{\theta_{J(C)}|_B}=B[2].
\]
Its type is $(2,2,2,2)$, and its kernel has order $2^8$. With $i$ acting by $\sigma$, the integral Hermitian lattice for $\xi$ is $U\oplus U$, where $U=\begin{psmallmatrix}0&1\\1&0\end{psmallmatrix}$.

The full integral calculation is HD-S01--HD-S02 below, applied at $g=3$: anti-invariant handle pairs have alternating pairing $2$, division by two gives a unimodular Riemann form, and the basis change $(A,B)\mapsto(A,iB)$ changes $\begin{psmallmatrix}0&i\\-i&0\end{psmallmatrix}$ into $U$. This also specifies the actual group scheme kernel, through the polarization homomorphism $\lambda_{2\xi}=\lambda_\xi\circ[2]$. The corresponding published identification is van Geemen \cite[Theorem~5.3(1)--(2), pp.~629--630]{vG96}.

\paragraph{The period rank and its exact locator.}
The codifferential for the genus-three cover, with its order-four line bundle $\eta$, is
\[
 H^0(X,\omega_X\eta)\otimes H^0(X,\omega_X\eta^{-1})
   \longrightarrow H^0(X,\omega_X^2).
\]
Its generic rank is exactly four: van Geemen \cite[Theorem~5.3(3) and its proof, pp.~629--630]{vG96}, extending Schoen \cite[Lemma~3.11 and the ensuing pencil calculation, pp.~29--30]{Schoen88}. This locates the Gaussian assertion in the published literature. Schoen's printed Theorem~3.2 is the Eisenstein statement.

For the requested explicit rank check, both source factors have dimension two. The base point free pencil sequence identifies the kernel with $H^0(X,\eta^{-2})=0$. Base point freeness holds at a general genus-three curve: a base point would give $\eta=\mathcal O(q-p)$, so $4q\sim4p$. The resulting degree-four map to $\mathbb P^1$ has two fully ramified values and at most six further branch values, by Riemann--Hurwitz. Its Hurwitz parameters have dimension at most $8-3=5$, whereas $\dim\mathcal M_3=6$. Thus the required multiplication rank is $2\cdot2=4$ on a dense open set. This paragraph checks the published proof in the precise order-four form used here; it is not a new dominance theorem.

\subsection{HD-Q02: the corrected Prym--B\'ezout calculation}
Schoen's canonical-linear-system supplier gives an effective detecting surface on $B$ of $\xi$-degree at most
\begin{equation}\label{eq:gaussian-prym-bound}
 D_{\rm Prym}=914\,248\,581\,120.
\end{equation}
The computation uses the full norm fibre as its support before applying $1-\sigma^2$.

\begin{proof}
Write $P=\ker\operatorname{Nm}_{C/X}^0$. This abelian variety has dimension six. Its restricted theta type is $(1,1,1,1,4,4)$. One can read the type directly in the handle description HD-S01: on each of the two regular handle pairs, the augmentation lattice has basis $e_0-e_1,e_1-e_2,e_2-e_3$, whose pairing matrix is the Cartan matrix of $A_3$ and has Smith invariants $(1,1,4)$. There are two such orthogonal blocks. Hence
\[
 \int_P\theta_{J(C)}^6=6!\,16.
\]
This establishes the elementary divisors, including their exponent, through an explicit integral pairing.

In Schoen's unramified construction, the inverse image of $|K_X|\simeq\mathbb P^2$ under $\operatorname{Sym}^4C\to\operatorname{Sym}^4X$ is a union of surfaces. At least one component gives a nonzero primitive determinant class after transfer; this is \cite[Theorem~2.0, the $r=0$ construction on pp.~12--13, and Corollary~3.1]{Schoen88}. Its Abel image lies in
\[
 W_4(C)\cap\operatorname{Nm}_{C/X}^{-1}([K_X]).
\]
Every connected component of the norm fibre is a translate of $P$. Each irreducible component of this intersection has dimension at most two: its preimage under the Abel map is contained in the inverse image of $|K_X|$, which is finite over $\mathbb P^2$.

We spell out why one of these surfaces detects after projection. On $C^4$, Schoen's nonzero primitive determinant projection belongs to
$U=(q\circ\operatorname{Abel})^*W_K(B)$, with $q=1-\sigma^2$. The restriction of the middle intersection pairing to $U$ is nondegenerate. Indeed $U$ lies in the tensor summand $H^1(C)^{\otimes4}$, which is killed by cup product with the product polarization (each summand would put degree three on a curve factor). Thus $U$ is primitive and of type $(2,2)$, and Hodge--Riemann makes its middle pairing positive definite. Consequently some effective supplier component pairs nontrivially with a class pulled back from $W_K(B)$. The projection formula forces both its Abel image $Z\subset J(C)$ and its image $q(Z)\subset B$ to have dimension two and nonzero pairing with a Weil class. Thus $Z$ is an irreducible component of the displayed intersection in a translate $T_0$ of $P$.

Embed $J(C)$ by $3\Theta$. The Poincar\'e formula and the computed Prym type give
\[
 \deg_{3\Theta}W_4=3^4\frac{9!}{5!}=244\,944,
 \qquad
 \deg_{3\Theta}T_0=3^6\,6!\,16=8\,398\,080.
\]
The component form of B\'ezout, including improper intersections, is Fulton \cite[Example~8.4.6]{Fulton98}. It gives
\[
 \deg_\Theta Z\le\frac{244\,944\cdot8\,398\,080}{3^2}
       =228\,562\,145\,280.
\]
The Rosati identity $q^\dagger q=2(1-\sigma^2)$ has eigenvalues zero and four. Therefore $q^*(\Theta|_B)\le4\Theta$ in the nef order. Since $\Theta|_B=2\xi$, it follows that $q^*\xi\le2\Theta$. Intersections with nef classes on the effective surface $Z$ give
\[
 \deg_\xi q_*Z=\int_Z(q^*\xi)^2\le4\deg_\Theta Z.
\]
The pushforward is effective integral and detects. Its reduced support contains a detector with no greater degree, proving \eqref{eq:gaussian-prym-bound}.
\end{proof}

The fourfold $B$ and the support $P$ are related by the actual morphism $q|_P:P\to B$; its connected kernel is the sign-character factor of dimension two. The norm condition in this construction therefore supplies a six-dimensional translate $T_0$. The proposed use of $3^4\,4!\,16$ for the degree of this $T_0$ is replaced by the six-dimensional value above. The polarization calculation on $B$ remains exactly $(2,2,2,2)$.

\subsection{HD-Q03: a degree-48 detector from the published theta construction}
On every fibre of the even principal Gaussian Weil-fourfold family with Hermitian lattice $U^2$,
\begin{equation}\label{eq:D0-Gaussian}
 \boxed{\delta_W\le D_0(\mathbb Q(i))=48.}
\end{equation}

\begin{proof}
We use the following precise published inputs from \cite{vG96}: Lemma~3.3, pp.~621--622, gives a finite base locus of sixteen points for the six odd sections of $2\Theta$ on an indecomposable member $X_\tau$; Proposition~3.6, pp.~623--624, places their rational map in a smooth quadric $Q\subset\mathbb P^5$; the proof of Theorem~3.7, pp.~624--625, proves dominance and a nonzero Weil projection for the strict inverse image of a general ruling plane $\Lambda\subset Q$ at a general member. Theorem~5.3(2) identifies this principal family with the lattice in HD-Q01.

Fix such a general member and one integral component $Y$ of that strict inverse image with $e_W[Y]\ne0$. Choose a general linear $\mathbb P^3$ containing $\Lambda$. Its intersection with $Q$ is the union of $\Lambda$ and a residual plane. Write the $\mathbb P^3$ as the intersection of two hyperplanes of $\mathbb P^5$. Their inverse images define two effective divisors $D_1,D_2\in|2\Theta|$ on $X_\tau$.

These divisors can be chosen with no common prime divisor. Indeed, the rational map is dominant and generically finite. On its domain of definition, a common prime divisor would be contracted to a subset of one of the two planes. There are finitely many contracted prime divisors: their closures on a resolution belong to the finitely many divisorial components of the exceptional locus. A general ruling plane contains none of their images, and the residual plane can be chosen to contain none either. The original base locus is finite and supplies no prime divisor. Thus $D_1\cdot D_2$ is an effective integral surface cycle, and $Y$ occurs in it with positive integral multiplicity.

Set $C=D_1\cdot D_2$ and $R=C-Y$. Then $R$ is effective integral,
\[
 [C]=4\theta^2,\qquad e_W[C]=0,\qquad
 e_W[R]=-e_W[Y]\ne0,
\]
and
\[
 \deg_\theta Y+\deg_\theta R
 =\deg_\theta C=4\int_{X_\tau}\theta^4=4\cdot4!=96.
\]
One of $Y,R$ consequently has degree at most $48$, and some integral component of that effective detecting cycle detects. This proves the bound at a general member. The closedness of $\{\delta_W\le48\}$, supplied as ATTACK Theorem D, extends it to the entire connected family, after finite level and trivialization of the Weil plane as in ATTACK Corollary W'.
\end{proof}

Combining \eqref{eq:D0-Gaussian} with HD-D02 gives the requested integral division constant
\begin{equation}\label{eq:C4-Gaussian}
 \boxed{C_4(\mathbb Q(i);\delta,T)
 =141\,557\,760\,\delta\,\lceil\sqrt T\rceil.}
\end{equation}
There are positive integers $N\le C_4$ and effective integral surfaces $Z^+,Z^-$ of $\theta$-degree at most $C_4$ with $Nt=[Z^+]-[Z^-]$ on every fibre. The class normalization and polarization in this statement are those fixed at the beginning of Section~\ref{sec:quadratic-extraction}.

\subsection{HD-Q04: the standard odd principal lattice, with its actual map}
Let $I_{1,1}=\operatorname{diag}(1,-1)$ and
$P=\begin{psmallmatrix}1&1\\1&-1\end{psmallmatrix}$. Direct multiplication gives
\[
 P^tUP=2I_{1,1},\qquad P^tI_{1,1}P=2U.
\]
On two blocks this is a $\mathbb Z[i]$-linear polarized isogeny between the corresponding even and odd period data, of degree $16$, with polarization pullback factor two. To see that this specifies an actual holomorphic morphism, transport the complex structure by this rational linear map and use its integral lattice inclusion; that construction is the defining period correspondence. The inverse rational linear map, multiplied by two, gives the reverse isogeny. Pushing forward a detecting surface multiplies its degree by at most $2^2=4$, and the induced isomorphism of rational Weil planes preserves nonzero projection. Hence the standard odd principal family $I_{2,2}$, containing the balanced product principally polarized $E_i^4$, satisfies
\[
 \delta_W\le192,\qquad
 c_t\le2\,264\,924\,160\,\delta\,\lceil\sqrt T\rceil.
\]
These two principal data are specified separately by their integral forms and by the displayed morphisms between them. No integral isomorphism between the two lattices is used.
